\chapter{Классификация вещественной формы коэффициентного цикла Тёплица}

Предыдущий гейт построил сбалансированную пару ориентированных операторов
Тёплица и доказал равенство нулю её обычного комплексного индекса. Однако
это ещё не отвечает на вопрос о вещественном классе: группа $KO_6(A)$
зависит не только от оператора, но и от конкретной вещественной формы
коэффициентной алгебры $A$.

Настоящий гейт не вводит нового оператора. Он сравнивает все естественные
вещественные формы уже построенной пары и выбирает ту, которая действительно
реализует проектный обмен
\begin{equation}
 L\longleftrightarrow L^*,\qquad
 E\longleftrightarrow E^*,\qquad
 T_+\longleftrightarrow T_-.
 \label{eq:v5-kr-project-exchange}
\end{equation}

\section{Литературная рамка}

Бурсема и Шокет систематически различают $KO$ вещественной алгебры,
комплексную $K$-теорию и $KR$ пространства с инволюцией; особенно важно,
что вещественная коммутативная алгебра не обязана иметь вид
$C(X,\mathbb R)$; см. \path{arXiv:2407.05880v4}. Шик строит точные
последовательности, связывающие вещественную теорию с теорией её
комплексификации; см. \path{arXiv:math/0311295}. Бурсема и Лоринг дают
представители всех восьми $KO$-групп унитариями с симметриями и явные
граничные отображения; см. \path{arXiv:1504.03284}. Следовательно,
сокращение двух комплексных индексов нельзя автоматически объявлять
обнулением вещественного класса: сначала требуется назвать вещественную
алгебру и отображение комплексификации.

\section{Три возможные вещественные формы}

Положим $B=M_{105}(\mathbb C)$. Для одной ориентированной ветви наиболее
очевидна поточечная инволюция
\begin{equation}
 \rho_{\rm pt}(a)=\overline a.
\end{equation}
Её неподвижная алгебра равна
\begin{equation}
 B^{\rho_{\rm pt}}=M_{105}(\mathbb R).
\end{equation}
По инвариантности Мориты
\begin{equation}
 KO_6\bigl(M_{105}(\mathbb R)\bigr)
 \simeq KO_6(\mathbb R)=0.
 \label{eq:v5-kr-pointwise-zero}
\end{equation}
Но эта форма сохраняет одну ветвь и не реализует обмен
\eqref{eq:v5-kr-project-exchange}. Поэтому она математически допустима,
но не является вещественной структурой текущего проектного цикла.

Вторая возможность --- кватернионная форма. Она несовместима с
$J^2=+1$ текущего KO6. Кроме того, комплексная размерность глобального
кватернионного модуля должна быть чётной, тогда как коэффициентный ранг
равен $105$. Эта ветвь исключается до индексного вычисления.

Наконец, рассмотрим обе ориентации совместно:
\begin{equation}
 A_{\rm ex}=B\oplus B,
 \qquad
 \rho_{\rm ex}(a,b)=(\overline b,\overline a).
 \label{eq:v5-kr-exchange-algebra}
\end{equation}
Её вещественная неподвижная алгебра имеет вид
\begin{equation}
 A_{\rm ex}^{\rho_{\rm ex}}
 =\{(a,\overline a):a\in B\}
 \simeq M_{105}(\mathbb C)_{\mathbb R},
 \label{eq:v5-kr-fixed-complex-real}
\end{equation}
где индекс $\mathbb R$ означает: комплексная матричная алгебра
рассматривается как вещественная $C^*$-алгебра.

Именно \eqref{eq:v5-kr-exchange-algebra} реализует действие проектного
$J$: оно не фиксирует ориентированные половины по отдельности, а меняет
их местами. Поэтому дальнейшая классификация должна производиться над
$A_{\rm ex}^{\rho_{\rm ex}}$, а не над $M_{105}(\mathbb R)$.

\section{Группа, допускающая целочисленный остаток}

Для комплексных чисел, рассматриваемых как вещественная алгебра,
вещественная $K$-теория совпадает с двухпериодической комплексной таблицей:
\begin{equation}
 KO_{2k}(\mathbb C_{\mathbb R})\simeq\mathbb Z,
 \qquad
 KO_{2k+1}(\mathbb C_{\mathbb R})=0.
\end{equation}
По инвариантности Мориты отсюда следует
\begin{equation}
 \boxed{
 KO_6\bigl(A_{\rm ex}^{\rho_{\rm ex}}\bigr)
 \simeq KO_6(\mathbb C_{\mathbb R})
 \simeq\mathbb Z.}
 \label{eq:v5-kr-selected-group}
\end{equation}

Это меняет прежний предварительный вывод. Нулевой индекс полного
комплексного оператора
\begin{equation}
 (-15)+(+15)=0
\end{equation}
не доказывает равенство нулю класса в группе
\eqref{eq:v5-kr-selected-group}. Комплексификация вещественной алгебры
$\mathbb C_{\mathbb R}$ раскладывается на две сопряжённые комплексные
компоненты, и противоположные ориентированные индексы могут быть двумя
тенями одного вещественного целочисленного класса.

Коэффициентный проектор пары имеет вид
\begin{equation}
 q_{\rm ex}=(q_0,\overline{q_0}),
 \qquad
 \operatorname{rank}_{\mathbb C}q_0=15.
\end{equation}
Поэтому естественный кандидат в
$KO_6(A_{\rm ex}^{\rho_{\rm ex}})\simeq\mathbb Z$ имеет абсолютную
величину пятнадцать. Его нормированное следовое чтение остаётся
\begin{equation}
 \frac{15}{105}=\frac17
 \qquad\text{или, на удвоенном носителе,}\qquad
 \frac{30}{210}=\frac17.
 \label{eq:v5-kr-one-seventh-survives}
\end{equation}

\begin{proposition}[Различие двух нулей]
Нулевой обычный индекс сбалансированного комплексного оператора является
суммой двух ориентированных компонент. Он не совпадает по смыслу с
утверждением, что вещественный класс в
$KO_6(M_{105}(\mathbb C)_{\mathbb R})$ равен нулю. Выбранная проектом
обменно-сопрягающая вещественная форма допускает целочисленный класс.
\end{proposition}

\section{Что ещё не доказано}

Формула \eqref{eq:v5-kr-selected-group} классифицирует возможную группу,
но ещё не вычисляет класс конкретного цикла. Для доказательства числа
$\pm15$ необходимо построить явный $Cl_{0,6}$-линейный фредгольмов
модуль над $A_{\rm ex}^{\rho_{\rm ex}}$ и проверить, что его отображение
в комплексную пару даёт именно ориентированные классы $-15$ и $+15$ с
требуемым правилом знаков.

\begin{nogocriterion}[Запрет подмены группы классом]
Из изоморфизма
\[
 KO_6(A_{\rm ex}^{\rho_{\rm ex}})\simeq\mathbb Z
\]
не следует автоматически, что построенная пара представляет число $15$.
До явной Clifford-линейной конструкции величина $15$ является
кандидатом, а не доказанным вещественным индексом.
\end{nogocriterion}

Тем более из этого результата пока не следуют локализация частицы,
конечная энергия, масса или динамический масштаб. Получен более узкий и
проверяемый математический вопрос: реализует ли уже построенная пара
ненулевой генератор выбранной целочисленной $KO_6$-группы?

\begin{protocol}
\begin{enumerate}
 \item поточечная вещественная форма одной ветви:
 \textbf{$M_{105}(\mathbb R)$, $KO_6=0$};
 \item соответствие этой формы проектному обмену $L\leftrightarrow L^*$:
 \textbf{нет};
 \item кватернионная форма: \textbf{исключена знаками и нечётным рангом};
 \item обменно-сопрягающая форма пары:
 \textbf{$M_{105}(\mathbb C)_{\mathbb R}$};
 \item выбранная группа: \textbf{$KO_6\simeq\mathbb Z$};
 \item сокращение комплексных индексов доказывает нулевой Real-класс:
 \textbf{нет};
 \item кандидат в абсолютный вещественный индекс: \textbf{$15$};
 \item нормированный вес: \textbf{$1/7$};
 \item явный $Cl_{0,6}$-линейный цикл: \textbf{не построен};
 \item класс $\pm15$ в $KO_6$: \textbf{не доказан};
 \item физическое замыкание: \textbf{не получено};
 \item следующий гейт:
 \textbf{явный $Cl_{0,6}$-линейный индекс обменного цикла}.
\end{enumerate}
\end{protocol}

Машинный сертификат сохранён в
\path{s2t_v5_real_toeplitz_kr_classification_gate_results.json}.